\section*{NeurIPS Paper Checklist}

\begin{enumerate}

\item {\bf Claims}
    \item[] Question: Do the main claims made in the abstract and introduction accurately reflect the paper's contributions and scope?
    \item[] Answer: \answerYes{}
    \item[] Justification: The abstract limits the strongest negative claim to the primary MLP backbone and the hardest Widar\_BVP setting, while explicitly noting that MLP+BN succeeds on SignFi. Up to 12 methods (Widar) and 10--11 (NTU-Fi/SignFi) are evaluated; the full coverage matrix is in Appendix~\ref{app:coverage}. Claims match the stated scope.

\item {\bf Limitations}
    \item[] Question: Does the paper discuss the limitations of the work performed by the authors?
    \item[] Answer: \answerYes{}
    \item[] Justification: Section~7 (Limitations) lists specific limitations including dataset scope, architecture, sample size, session splits, method fidelity, source overfitting, missing channel estimation, claim scope, and release/provenance constraints; an extended reviewer-FAQ is provided in Appendix~\ref{app:reviewer_faq}.

\item {\bf Theory assumptions and proofs}
    \item[] Answer: \answerNA{}
    \item[] Justification: The paper is a benchmark and empirical study. The root cause analysis is evidence-based, not theorem-proof style.

\item {\bf Experimental result reproducibility}
    \item[] Answer: \answerYes{}
    \item[] Justification: Appendix provides full hyperparameter table, reproduction commands, pinned environment (\texttt{requirements-lock.txt}), split-manifest description, and 248 automated tests. The anonymised repository at \url{https://anonymous.4open.science/r/wifi-tta-bench-ed2026} provides a single-command reproducer (\texttt{python scripts/reproduce\_all\_results.py --force}) plus fail-fast per-component scripts whose outputs back every table and figure in the paper (see \texttt{manuscript/paper2/SUBMISSION\_MANIFEST.md}).

\item {\bf Open access to data and code}
    \item[] Answer: \answerYes{}
    \item[] Justification: Source code is accessible at the double-blind anonymised URL \url{https://anonymous.4open.science/r/wifi-tta-bench-ed2026}. The Croissant 1.0 (core + RAI) dataset card is mirrored at \url{https://huggingface.co/datasets/anonymous-ed2026/wifi-tta-bench}; it is also included in the supplementary bundle at \texttt{outputs/croissant/wifi\_tta\_bench\_metadata.json}. Raw Widar\_BVP, NTU-Fi, and SignFi-10 CSI recordings are \emph{not} redistributed and remain under their upstream academic licences; WiFi-TTA-Bench ships split manifests, preparation scripts, evaluation code, and canonical per-method result JSONs. No redistributed benchmark artifact exceeds the 4\,GB sample-file threshold, so the reviewer-visible artifacts are directly inspectable without personal request.

\item {\bf Experimental setting/details}
    \item[] Answer: \answerYes{}
    \item[] Justification: Sections~3--4 specify all datasets (Table~1), methods and fidelity labels (Table~2), physics-guided variants, backbone architecture, hyperparameters, and evaluation protocol with 5 metrics.

\item {\bf Experiment statistical significance}
    \item[] Answer: \answerYes{}
    \item[] Justification: All results report 95\% bootstrap CIs, Cohen's $d$, and negative adaptation rate. With 11 TTA methods compared against \texttt{no\_adapt}, Bonferroni-corrected threshold is $\alpha_{\text{corr}} \approx 0.0045$; readers are advised to apply this when interpreting CIs.

\item {\bf Experiments compute resources}
    \item[] Answer: \answerYes{}
    \item[] Justification: Full benchmark runs in approximately 30 minutes on CPU. No GPU required. Python 3.12, PyTorch 2.10.

\item {\bf Code of ethics}
    \item[] Answer: \answerYes{}
    \item[] Justification: No new human subjects data are collected. Privacy implications of WiFi sensing and acceptable-use guidance are discussed in Section~7 and Appendix~\ref{app:reviewer_faq}.

\item {\bf Broader impacts}
    \item[] Answer: \answerYes{}
    \item[] Justification: Section~7 and Appendix~\ref{app:reviewer_faq} discuss that WiFi CSI sensing raises privacy concerns and that improved adaptation could increase surveillance capability.

\item {\bf Safeguards}
    \item[] Answer: \answerNA{}
    \item[] Justification: The benchmark does not release scraped data and does not redistribute raw datasets with restricted licenses. The release includes acceptable-use guidance against covert monitoring.

\item {\bf Licenses for existing assets}
    \item[] Answer: \answerYes{}
    \item[] Justification: Widar3.0 is used under its academic license. NTU-Fi is used under academic terms. Licenses noted in the datasheet.

\item {\bf New assets}
    \item[] Answer: \answerYes{}
    \item[] Justification: WiFi-TTA-Bench is a new benchmark asset released under the NeurIPS 2026 Evaluations \& Datasets track. Machine-readable metadata in Croissant 1.0 format with core and Responsible-AI (RAI) fields is provided at \texttt{outputs/croissant/wifi\_tta\_bench\_metadata.json}. Documentation also includes dataset cards (Appendix~\ref{app:datasheet}), split manifests, evaluation protocol (Appendix~\ref{app:hyperparams}), and acceptable-use guidance (Appendix~\ref{app:reviewer_faq}). Raw upstream datasets remain under their original academic licences; the benchmark redistributes only evaluation artifacts and split manifests.

\item {\bf Crowdsourcing and research with human subjects}
    \item[] Answer: \answerNA{}
    \item[] Justification: No crowdsourcing or human subjects research.

\item {\bf Institutional review board (IRB) approvals}
    \item[] Answer: \answerNA{}
    \item[] Justification: No human subjects research.

\item {\bf Declaration of LLM usage}
    \item[] Answer: \answerYes{}
    \item[] Justification: Claude (Anthropic) was used as a coding assistant for implementation, testing, and draft editing. It is not a component of any benchmark method. Full disclosure is in Appendix~\ref{app:ai}.

\end{enumerate}
